% ---------------------------------------------------------------------------
% Phase 6 — M01b3 (★ winner) — Style-content mixing on FLUX double-stream.
%
% Drop-in snippet for the experiments chapter (e.g. chapter 4, section
% "Смешивание стиля и содержания через B-LoRA"). Requires booktabs + tabularx.
%
% Figure asset path: figures/phase6_m01b3_grid.png
%   (origin: s3://.../20260515T112125/m01_mixing_grid_v3_with_suffix/grid.png)
% Source data:      output/results/m01_mixing_grid_v3_with_suffix_metrics.json
% ---------------------------------------------------------------------------

\subsection{Качественная оценка матрицы $3 \times 3$}
\label{subsec:phase6_m01b3_qualitative}

\begin{figure}[ht]
    \centering
    \includegraphics[width=\textwidth]{figures/phase6_m01b3_grid.png}
    \caption{Сетка результатов смешивания B-LoRA на FLUX.1-dev (этап M01b3): по строкам --- стилевые источники (\textit{Красные виноградники в Арле}, 1888; \textit{Звёздная ночь}, 1889; \textit{Стог сена в Живерни}, 1886), по столбцам --- контентные запросы (\textit{кот на стуле}, \textit{бегущая собака}, \textit{рюкзак на столе}). В каждой ячейке применено блочное слияние LoRA-адаптеров: стилевые ключи для \texttt{transformer\_blocks.0--8} и контентные --- для \texttt{transformer\_blocks.9--18}, с общим масштабом \texttt{lora\_scale}\,=\,0{,}7; для активации стилевого канала к промпту добавлен суффикс \textit{<<painted in the style of \dots>>}. В 6 из 9 ячеек наблюдается визуально различимый перенос стиля при сохранении исходного содержания, что эмпирически подтверждает гипотезу о разделимости стиля и содержания в блочной структуре FLUX. Ячейка (строка 1, столбец 3) демонстрирует утечку запоминания: исходное полотно воспроизводится как плакат на стене, что задокументировано далее как известное ограничение LoRA на одной обучающей картине.}
    \label{fig:phase6_m01b3_grid}
\end{figure}

Как видно из \autoref{fig:phase6_m01b3_grid}, предложенная блочная декомпозиция B-LoRA эмпирически переносится на трансформерную архитектуру FLUX: в 6 из 9 ячеек сетки стилевой канал \texttt{transformer\_blocks.0--8} формирует мазок и палитру источника, тогда как контентный канал \texttt{transformer\_blocks.9--18} удерживает геометрию запроса --- асимметрия средних значений DINO-content и DINO-style составляет около $5{,}7\times$ ($0{,}562$ против $0{,}099$). Существенно, что эффект разделимости проявляется лишь \emph{при условии} текстовой активации стиля через суффикс промпта; без него (конфигурации v1/v2) аналогичное блочное слияние не давало визуально различимого переноса, что мы трактуем как методологическое ограничение протокола. Отдельно следует отметить аномалию в ячейке (\textit{Красные виноградники}, \textit{рюкзак}) со значением DINO-style $= 0{,}275$: вместо стилизации модель воспроизводит обучающее полотно дословно в виде постера на стене, что иллюстрирует известный режим отказа LoRA --- запоминание единичного обучающего примера.

\subsection{Количественная оценка матрицы $3 \times 3$}
\label{subsec:phase6_m01b3_quantitative}

Для верификации качественных наблюдений, полученных при визуальном анализе сетки смешивания M01b3, проведена количественная оценка девяти ячеек по двум независимым осям с использованием самосупервизионного бэкбона DINO ViT-B/8. Протокол измерения следующий: для каждой ячейки матрицы $3 \times 3$ вычисляется (i) косинусное сходство между её DINO-эмбеддингом и эмбеддингом одиночного изображения-эталона стиля соответствующей строки --- метрика \textit{DINO-style} --- и (ii) косинусное сходство с усреднённым эмбеддингом по всему тренировочному набору изображений содержания соответствующего столбца --- метрика \textit{DINO-content}. Выбор DINO в качестве экстрактора признаков мотивирован его доказанной устойчивостью к низкоуровневым искажениям и способностью кодировать одновременно семантику объекта и стилистические характеристики при разной глубине проекции.

Полные значения покомпонентных метрик, а также средние по строкам и столбцам с большим средним по матрице, приведены в~\autoref{tab:phase6_mixing_dino} (Панель~A для стиля, Панель~B для содержания). Все генерации выполнены с активирующим стиль суффиксом в текстовом запросе на основании установленной в Phase~6 закономерности: без явного словесного вызова стиля LoRA-веса на FLUX не активируются и обнуляют перенос стиля во всех девяти ячейках.

\begin{table}[ht]
\centering
\caption{Количественная оценка disentanglement на матрице $3 \times 3$: косинусные сходства DINO ViT-B/8 между сгенерированными ячейками и эталонами стиля / содержания.}
\label{tab:phase6_mixing_dino}
\small
\begin{tabularx}{\textwidth}{l *{3}{Y} Y}
\toprule
\multicolumn{5}{c}{\textbf{Панель A. DINO-style} (cos с одиночным изображением-эталоном стиля строки), $\uparrow$ --- лучше} \\
\midrule
\textbf{Стиль (строки) $\backslash$ Содержание (столбцы)} & \textbf{cat} & \textbf{dog} & \textbf{backpack} & \textbf{Среднее по строке} \\
\midrule
van\_gogh\_starry\_road  & \textbf{0,115} & 0,045 & 0,275\textsuperscript{\dag} & 0,145 \\
van\_gogh\_starry\_night & \textbf{0,082} & 0,037 & 0,080 & 0,066 \\
monet\_garden            & 0,054 & \textbf{0,163} & 0,041 & 0,086 \\
\midrule
\textbf{Среднее по столбцу} & 0,083 & 0,082 & 0,132 & \textbf{0,099} \\
\bottomrule
\end{tabularx}

\vspace{0.6em}

\begin{tabularx}{\textwidth}{l *{3}{Y} Y}
\toprule
\multicolumn{5}{c}{\textbf{Панель B. DINO-content} (cos со средним эмбеддингом тренировочной папки содержания), $\uparrow$ --- лучше} \\
\midrule
\textbf{Стиль (строки) $\backslash$ Содержание (столбцы)} & \textbf{cat} & \textbf{dog} & \textbf{backpack} & \textbf{Среднее по строке} \\
\midrule
van\_gogh\_starry\_road  & 0,609 & 0,598 & 0,449 & 0,552 \\
van\_gogh\_starry\_night & 0,623 & 0,763 & 0,485 & 0,624 \\
monet\_garden            & 0,589 & 0,447 & 0,497 & 0,511 \\
\midrule
\textbf{Среднее по столбцу} & 0,607 & 0,602 & 0,477 & \textbf{0,562} \\
\bottomrule
\end{tabularx}

\vspace{0.4em}
\footnotesize
\textsuperscript{\dag} Ячейка (van\_gogh\_starry\_road, backpack) демонстрирует утечку памяти LoRA: модель воспроизводит фрагмент тренировочной картины <<Звёздная ночь>> в виде постера, что искусственно завышает DINO-style и не должно интерпретироваться как успешный перенос стиля. Полужирным в Панели~A выделены ячейки $\arg\max$ по строке среди \emph{валидных} (не подверженных утечке).
\end{table}

\normalsize

Большое среднее DINO-content $= 0{,}562$ превосходит большое среднее DINO-style $= 0{,}099$ примерно в $5{,}7$ раза. Подобная асимметрия является косвенным, но систематическим свидетельством того, что content-блоки FLUX в схеме B-LoRA доминируют над style-блоками: смысл объекта сохраняется на эталонном уровне, тогда как стилистический сдвиг остаётся умеренным даже после prompt-side активации. Это согласуется с предположением о том, что DiT-архитектура FLUX концентрирует семантическое содержание в небольшом числе блоков, легко покрываемых низкоранговой адаптацией, тогда как стиль распределён более диффузно и требует либо большего ранга, либо дополнительного prompt-уровневого вмешательства.

На уровне отдельных строк наблюдается выраженная неоднородность переноса стиля: $\overline{\mathrm{DINO\text{-}style}}_{\mathrm{vg\_road}} = 0{,}145$, $\overline{\mathrm{DINO\text{-}style}}_{\mathrm{monet}} = 0{,}086$, $\overline{\mathrm{DINO\text{-}style}}_{\mathrm{vg\_night}} = 0{,}066$. Парадоксально низкий результат строки \textit{van\_gogh\_starry\_night}, формально наиболее <<узнаваемого>> эталона, количественно подтверждает качественное наблюдение о том, что во всех трёх ячейках этой строки стиль фактически не активируется и изображения остаются фотореалистичными. Возможной причиной является сильная семантическая перегрузка токена <<starry night>> в текстовом энкодере, конкурирующая с действием LoRA-весов и приводящая к тому, что условный сигнал стиля гасится prompt-альтернативой того же концепта.

Ячейка (van\_gogh\_starry\_road, backpack) с DINO-style $= 0{,}275$ формально является глобальным максимумом стиля по всей матрице, однако соответствующее значение DINO-content $= 0{,}449$ --- минимум по строке. Такое смещение количественно подтверждает зафиксированную при визуальном анализе утечку памяти LoRA: модель буквально воспроизводит <<Звёздную ночь>> как постер вместо переноса её стилистики на рюкзак. Высокое DINO-style объясняется не успешным disentanglement, а копированием пикселей тренировочного изображения; именно поэтому в Панели~A таблицы данная ячейка помечена сноской и исключена из выделения полужирным.

Содержание демонстрирует существенно более стабильное поведение, однако и здесь столбец \textit{backpack} имеет систематически наименьшее среднее DINO-content $= 0{,}477$ против $0{,}607$ для \textit{cat} и $0{,}602$ для \textit{dog}. Вероятное объяснение состоит в том, что DINO ViT-B/8, обученный преимущественно на ImageNet-подобной семантике, лучше кластеризует одушевлённые категории с богатой внутриклассовой формой, тогда как ригидный искусственный объект <<рюкзак>> представлен в эмбеддинг-пространстве более узким кластером и допускает меньший геометрический и текстурный дрейф без потери близости. Дополнительным фактором служит специфическая чувствительность рюкзака к утечке стиля (см. ячейку (0, backpack)), которая вытесняет содержательный сигнал.

\subsubsection*{Скалярные показатели disentanglement}

Для интегральной оценки баланса осей вводятся три скалярных показателя.

\paragraph{Style-content tradeoff index (STI).}
\begin{equation}
\mathrm{STI} = \overline{\mathrm{DINO\text{-}content}} - \overline{\mathrm{DINO\text{-}style}} = 0{,}562 - 0{,}099 = 0{,}463.
\label{eq:phase6_sti}
\end{equation}
Положительное и большое по модулю значение $\mathrm{STI} \approx 0{,}463$ количественно фиксирует асимметрию системы в сторону сохранения содержания: контентная ось <<выигрывает>> у стилистической приблизительно на половину диапазона косинусного сходства, что характеризует текущую конфигурацию B-LoRA на FLUX как content-dominant.

\paragraph{Внутристрочная вариативность стиля (style row consistency).}
\begin{equation}
\overline{\sigma}_{\mathrm{style\_row}} = \mathrm{mean}(0{,}0963,\, 0{,}0206,\, 0{,}0546) = 0{,}0572.
\label{eq:phase6_style_row_std}
\end{equation}
Поскольку средний внутристрочный разброс DINO-style ($0{,}057$) сопоставим по порядку величины с самим средним значением DINO-style ($0{,}099$), стиль внутри строки нельзя признать инвариантным к содержанию --- степень стилистического переноса значимо зависит от объекта в ячейке, что является дополнительным симптомом неполного disentanglement.

\paragraph{Стабильность содержания по столбцам (content column stability).}
\begin{equation}
\overline{\sigma}_{\mathrm{content\_col}} = \mathrm{mean}(0{,}0137,\, 0{,}1291,\, 0{,}0207) = 0{,}0545.
\label{eq:phase6_content_col_std}
\end{equation}
Низкий разброс в столбцах \textit{cat} ($0{,}014$) и \textit{backpack} ($0{,}021$) свидетельствует о высокой устойчивости содержания к смене стиля; аномально высокий разброс столбца \textit{dog} ($0{,}129$) объясняется единственной ячейкой (monet, dog) с DINO-content $= 0{,}447$ против $\approx 0{,}68$ в двух других строках --- ровно той ячейкой, которая качественно признана <<самой чистой>> по стилю, что указывает на локальный обмен между осями именно в этой точке матрицы.

\subsubsection*{Итог Phase 6}

Совокупность приведённых количественных данных подтверждает центральный тезис работы в его уточнённой формулировке: disentanglement содержания и стиля посредством B-LoRA на FLUX реализуется, однако \emph{условно} --- при обязательной prompt-уровневой активации стиля и при отсутствии переобучения LoRA на единичную композицию-источник. Без этих условий перенос стиля либо подавляется текстовым энкодером (строка \textit{van\_gogh\_starry\_night}), либо вырождается в воспроизведение тренировочного изображения (ячейка (van\_gogh\_starry\_road, backpack)).
